\scanpage{146}
\mspage{74}
% archive sheet number: 145
% L1–7
et comme $G_0^+$ est transitif sur $S$, il s'ensuit bien que (43) est indépendant du choix de $s\in S$. [Si on ne supposait pas les $P_i$ réguliers, il y aurait lieu soit de mettre un axiome d'indépendance, soit de se fatiguer pour prouver cette indépendance par voie combinatoire, en termes des autres propriétés\ldots] Ceci dit, si $(\lambda,\omega,\omega')\in\Lambda_n\times\mu(i)\times\mu(\dot\lambda.i)$ se correspondent par (44) (pour un $s$, donc pour $\forall\,s$) on posera par définition
% L11–12
\[
\left\{\begin{aligned}
 \omega'&=\lambda.\omega &&(\omega\in\mu(i),\ \omega'\in\mu(i'\unsure{)|\vphantom{)}}=\mu(\dot\lambda.i)\\
 \Longleftrightarrow\quad\omega'_E&=\omega_E^\lambda
 \quad\substack{\text{pour tte étoile}\\\text{(ou une étoile)}}&&\lambda\in\Lambda_n^*).
\end{aligned}\right.
\tag*{(46)}
\]
% L12 left
% slanted
\margin{Il est clair déjà que $\forall\,\omega,\omega'\in\unsure{\vec\Pi|\vec\pi}$, $\exists1\ \lambda\in Z_{\deleted{\unsure{m}}n}=\Lambda_n^*$, avec $\omega'=\lambda.\omega$}
% L13–14
Soit de même $\lambda'\in\Lambda_n^*$, $i''=(\lambda'\lambda\deleted{\unsure{)}})i=\lambda'i'$, considérons $\omega''=\lambda'.\omega'\inserted{\in\mu(i'')}$, caractérisé par
% L15
\[
\omega''_E=(\omega'_E)^{\lambda'}\quad\text{donc}\quad
 \omega''_E=((\omega^{\unsure{*}}_E)^\lambda)^{\lambda'}=(\omega_E)^{\lambda\lambda'}
\]
% L16–17
on a donc aussi, par définition de $(\lambda\lambda').\omega^{*}$, $\omega''=(\lambda\lambda').\omega$ \deleted{\unsure{d'où}}, d'où la relation d'associativité
% L18
\[
\lambda'(\lambda\omega)=(\lambda'\lambda).\omega\,,\tag*{(47)}
\]
% L19–20
donc la loi (46) définit \deleted{bien} une opération du groupe $\Lambda_n^*$ sur $\vec\Pi$, \deleted{que} de sorte que
% L21
\[
\vec\Pi\quad\text{est un torseur sous }\Lambda_n^*\text{ via }(46).\tag*{(48)}
\]
% slanted
\margin{Ceci marche également pour le cas de l'octaèdre !}

% L22
\deleted{\unsure{du groupe|Le groupe}} Essayons de prouver que les conditions
